% Generated draft: requires mathematical and source review.
\documentclass[11pt]{article}
\usepackage[margin=25mm]{geometry}
\usepackage{fontspec}
\setmainfont{Latin Modern Roman}
\usepackage{amsmath,amssymb,mathtools,unicode-math}
\setmathfont{Latin Modern Math}
\usepackage{enumitem,needspace,xurl,hyperref}
\setlength{\parindent}{0pt}
\setlength{\parskip}{5pt}
\setlength{\emergencystretch}{3em}
\newcommand{\N}{\mathbb{N}}
\newcommand{\Z}{\mathbb{Z}}
\newcommand{\Q}{\mathbb{Q}}
\newcommand{\R}{\mathbb{R}}
\newcommand{\C}{\mathbb{C}}
\newcommand{\F}{\mathbb{F}}
\newcommand{\A}{\mathbb{A}}
\newcommand{\PP}{\mathbb{P}}
\newcommand{\E}{\mathbb{E}}
\newcommand{\eps}{\varepsilon}
\newcommand{\dd}{\,\mathrm{d}}
\begin{document}
\section*{UnsolvedMath ID: 30002355}
\textbf{Source problem code:} catalogue entry
\subsection{Definitions and mathematical statement}
Let \(X\) be an integral scheme such that \(2\) is invertible in \(\mathcal O_X\), and let \(K=K(X)\) be its function field. In particular, \(\operatorname{char}(K)\ne 2\). Let \(\Omega(K)\) denote the set of equivalence classes of all nontrivial rank-one discrete valuations of \(K\). A representative may be normalized as a surjective valuation \(v:K^\times\to\Z\); write \(K_v\) for the corresponding completion. As stated in the supplied record, this set comprises all such valuations, with no restriction to valuations having a center on \(X\).

Let \(V\) be a finite-dimensional \(K\)-vector space of positive dimension, and let \(q:V\to K\) be a quadratic form. The form is nondegenerate if its polar form \(b_q(x,y)=q(x+y)-q(x)-q(y)\) has zero radical. For a field extension \(L/K\), write \(q_L\) for the scalar extension to \(V\otimes_K L\). The form \(q_L\) is isotropic if there exists \(0\ne x\in V\otimes_K L\) with \(q_L(x)=0\).

The general question in the supplied record asks whether the following implication holds for every nondegenerate quadratic form \(q\) over \(K\):
\[
\left(\forall v\in\Omega(K),\ \exists\,0\ne x_v\in V\otimes_K K_v\text{ such that }q_{K_v}(x_v)=0\right)
\quad\Longrightarrow\quad
\left(\exists\,0\ne x\in V\text{ such that }q(x)=0\right).
\]
The supplied background reports counterexamples in some classes of fields. Thus this broad question is retained as a question about the scope of the principle, without asserting that the unrestricted implication is an unresolved conjecture.

The more specific question is the following. Let \(k\) be a totally imaginary number field, meaning a finite extension of \(\Q\) with no embedding into \(\R\), let \(C\) be an integral algebraic curve over \(k\), and put \(K=k(C)\). For every nondegenerate quadratic form \(q\) over \(K\) of dimension \(n\ge 5\), does the displayed implication hold with \(\Omega(K)\) consisting of all rank-one discrete valuations of \(K\)? The supplied record does not specify whether the curve must be geometrically integral, smooth, or proper, or whether additional conditions on its constant field are intended. These qualifications remain unresolved; none is imposed here.
\subsection{Short English statement}
The question asks when a quadratic form that has a nonzero zero over every discrete valuation completion also has a nonzero zero over the original field. Its specific target concerns forms of dimension at least \(5\) over function fields of curves over totally imaginary number fields. The precise intended hypotheses on the curve require review.
\subsection{Sources}
\begin{itemize}
\item Local UnsolvedMath registry, permanent ID \(30002355\), title \emph{Local--Global Isotropy over Function Fields}, source problem code \texttt{catalogue entry}, in the Oberwolfach Reports Open Problems collection. The record supplies proposed year \(2013\) and names no individual proposer.
\item Supplied report citation: \emph{Quadratic Forms and Linear Algebraic Groups} (2014), Oberwolfach Reports, pp.~1819--1859, DOI \texttt{10.4171/owr/2013/31}. \url{https://doi.org/10.4171/owr/2013/31}. The supplied background describes these reports as records of mathematical workshops and their open-problem discussions. The citation's year \(2014\), the DOI's \(2013\) component, and the proposed year \(2013\) are retained as supplied.
\item Supplied literature source attributed to Auel, Parimala and Suresh, labeled ``local-global principle'': \url{https://math.dartmouth.edu/~auel/papers/docs/AUEL_MFO2013.pdf}.
\item Supplied literature source attributed to Auel and Suresh, labeled ``quadratic-form counterexamples'': \url{https://doi.org/10.2140/ant.2024.18.1497}.
\item The literature assessment dated \texttt{2026-08-21} is supplied dated curation, not independently verified here. It labels the problem partially solved and reports that the principle holds in global fields and several low-dimensional cases, that counterexamples occur over surface and higher-transcendence-degree fields, and that no general classification for a given function field is known. The background explicitly requests independent verification of time-sensitive status and cited claims.
\item The registry labels the statement \texttt{corrected\_verified} and reports that it was checked against the official Oberwolfach report to restore the source hypotheses and notation. This is a supplied verification claim, not an independent verification performed for this draft.
\end{itemize}
\end{document}
